
\subsection{Simultaneous local strict covariance and nonzero slice torque}
\label{subsec:two-mode-torque-XXIV}
The preceding two-frequency field has zero mean angular torque in
the special choice $\lambda=0$. A small additional first-harmonic
polarization produces an explicitly nonzero torque without losing
the positive covariance floor.

\begin{proposition}[Divergence-free two-harmonic torque shaper]
\label{prop:two-mode-torque-XXIV}
Fix real $f,g\in C_c^\infty((1/2,3/2))$ such that
$f=g=1$ on $[3/4,5/4]$. Let $|\lambda|\le1/16$.
Use the $m=+1$ cylindrical amplitudes
\[
a=f(r)\cos z,\quad
b=ia+\lambda r^2f(r)\cos z,\quad
c=-f'(r)\sin z-i\lambda rf(r)\sin z
\]
and the independent $m=+2$ amplitude
$(a_2,b_2,c_2)=(0,0,g(r))$; add complex-conjugate modes
to make the field real. On $\mathbb R^2\times\mathbb S^1$
of axial period $2\pi$, the resulting smooth finite-energy
velocity is divergence-free. Its co-rotating azimuthal stress
$Q=\langle v\otimes v\rangle_\theta$ satisfies
\[
\boxed{\det Q=8 f(r)^4g(r)^2\cos^4z.}
\]
On the plateau
$3/4\le r\le5/4,\ |z|\le\pi/6$,
\begin{equation}\label{eq:strict-positive-torque-XXIV}
\boxed{\quad Q\succeq\frac{653}{512}I_3
             \succ I_3.\quad}
\end{equation}
Moreover, the radial-slice angular flux and torque are
\[
\boxed{
T(z)=\lambda I_f\sin 2z,\qquad
m(z)=2\lambda I_f\cos2z,\qquad
I_f=\int_0^\infty r^3 f(r)^2\,dr>0.
}
\]
In particular, the torque is nonzero when $\lambda\ne0$.
\end{proposition}
\begin{proof}
The first-harmonic divergence vanishes because
$a_r+a/r+ib/r+c_z=0$; the second mode has axial
velocity independent of $z$, and therefore also has zero
divergence. Compact annular support gives Cartesian
smoothness and finite energy.

Writing the real and imaginary first-harmonic polarizations as
\[
p=(f\cos z,\,\lambda r^2 f\cos z,\,-f'\sin z),\qquad
q=(0,\,f\cos z,\,-\lambda r f\sin z)
\]
and the independent second-harmonic polarization as
$s=(0,0,g)$ gives $Q=2(pp^{\mathsf T}+qq^{\mathsf T}
+ss^{\mathsf T})$. Since the determinant of the
three-column matrix $(p,q,s)$ is $f^2g\cos^2z$,
$\det Q=8f^4g^2\cos^4z$, independently of $\lambda$.

On the plateau the diagonal entries satisfy
$Q_{11}=2\cos^2z,\ Q_{22}\ge2\cos^2z,\ Q_{33}\ge2$.
The only nonzero off-diagonal terms are
$Q_{12}=2\lambda r^2\cos^2z$ and
$Q_{23}=-2\lambda r\sin z\cos z$.
Using $2|xy|\le x^2+y^2$ in the quadratic form gives a
uniform lower bound equal to the smallest of
\[
2\cos^2z(1-|\lambda|r^2),\quad
2\cos^2z(1-|\lambda|r^2)-2|\lambda|r|\sin z\cos z|,
\quad
2-2|\lambda|r|\sin z\cos z|.
\]
The hypotheses imply these are respectively bounded below
by $693/512$, $653/512$ and $123/64$, giving the stated
positive floor.

Angular Fourier orthogonality removes all mixed $m=1,m=2$
terms in $Q_{\theta z}$. The first pair contributes
\[
Q_{\theta z}=-\lambda r(rff'+f^2)\sin 2z.
\]
Integrating $r^2Q_{\theta z}$ in $r$ and using
$\int r^4ff'=-2\int r^3f^2$ proves
$T=\lambda I_f\sin2z$; differentiating yields the torque.
\end{proof}

\paragraph{Release boundary.}
The construction simultaneously passes three \emph{local}
kinematic gates---solenoidality, strict covariance positivity
and nonzero slice torque---but remains a manufactured field,
not a solution of the time-dependent unforced Navier--Stokes
equations. In particular it is not certified to lie in the
specific H8 seven-ray positive-coefficient cone, is not matched
to the A.11/A.19 source, and does not cancel $m=3,4$
quadratic mode leakage. Strict local covariance cannot
replace the global anti-divergence, pressure, source and
singularity-preservation proof obligations.
Independent symbolic determinant, full modal divergence,
angular torque and rational covariance-floor calculations pass.
